\documentclass[10pt,a4paper]{article}
\usepackage[margin=23mm]{geometry}
\usepackage[T1]{fontenc}
\usepackage{lmodern}
\usepackage{microtype}
\usepackage{amsmath,amssymb,booktabs,array}
\usepackage{xcolor}
\usepackage{listings}
\usepackage{enumitem}
\usepackage[colorlinks=true,urlcolor=blue!55!black,citecolor=blue!55!black,linkcolor=blue!55!black]{hyperref}
\hypersetup{pdftitle={Bring Your Own Body},pdfauthor={Octavia Sulea},pdfsubject={Lean2LLM: rechecked proof search and portable proof storage for Lean}}
\usepackage{fancyhdr}
\pagestyle{fancy}
\fancyhf{}
\fancyhead[L]{\small Lean2LLM}
\fancyhead[R]{\small Research prototype, 22 October 2026}
\fancyfoot[C]{\thepage}
\setlength{\headheight}{14pt}
\setlength{\parskip}{3pt}
\setlength{\parindent}{1em}
\setlist{nosep,leftmargin=*}
\lstset{basicstyle=\small\ttfamily,columns=fullflexible,breaklines=true,
  frame=single,rulecolor=\color{black!15},backgroundcolor=\color{black!2},
  showstringspaces=false,keepspaces=true}
\newcommand{\code}[1]{\texttt{#1}}
\title{Bring Your Own Body}
\author{Octavia Sulea\\\small\href{mailto:octavia@octav.ai}{octavia@octav.ai}\\
  \small\url{https://github.com/mary-octavia/lean-llm}}
\date{22 October 2026\\\small Software and systems manuscript; draft for author review}

\begin{document}
\maketitle
\begin{abstract}
Language models can propose formal proofs, while proof assistants check the
resulting proof terms. A useful research workflow must also preserve the exact
statement being proved, make verification conditions inspectable, and retain
artifacts that another user can find and check. Lean2LLM develops these concerns
around an existing local Lean--LLM bridge. Its proposal loop fixes the context
and proposition, supplies Lean's actual goal and diagnostics to an external
proposer, and accepts an exported proof only after Lean checking and a
transitive axiom audit. An exact-request cache rechecks every hit. A separate
content-addressed proof library adds indexed retrieval, explicit toolchain
identities, portable bundles, and local verification records. A browser
workbench checks restricted core Lean source in a fresh WebAssembly worker
and exchanges records with the Python library. Retrieval proposes
an artifact for use; a fresh check establishes its status in the receiving
environment. A second verification profile accommodates standalone kernel
certificates, drawing on a published Navier--Stokes reproducibility bundle.
This paper documents the architecture, artifact schema, trust boundaries, and
reproduction procedure. It reports inherited repository evidence separately
from tests of the new storage implementation. The contribution is a practical
composition of proof generation, checking, storage, and retrieval.
\end{abstract}

\section{A proof should remain usable after the chat}
An interactive proof experiment often ends with a successful compiler message
and a source file. Reusing that result later requires more information: the
fixed proposition, definitions in scope, proof body, permitted assumptions,
Lean version, dependency identity, and the procedure that produced the
verification result. A plausible explanation or a remembered success message
does little to resolve a mismatch in any of these fields.

Lean2LLM joins a bounded proof-proposal loop to an artifact library. The starting
point is the public \code{lean-llm} repository, whose native Lean tactic and Python
driver already implement proposal, feedback, kernel checking, export, and exact
cache replay~\cite{sulea-bridge}. The new library makes the stored proof an
explicit exchange object. A web interface exposes the same distinction between
draft source, a recorded verification event, and a check performed in the current
environment. The public-facing project name is Lean2LLM; the associated domains
are \url{https://lean2llm.com}, \url{https://lean4llm.com}, and
\url{https://langlean.si}. Domain names identify interfaces, while repository
revisions and artifact digests identify reproducible material.

There are three engineering objectives. First, the model must solve the given
goal under its fixed statement. Second, retrieval must make an
existing proof easy to discover without treating textual similarity as logical
equivalence. Third, a reader should be able to download enough evidence to
understand and repeat the relevant check. These objectives lead to two
verification profiles and two complementary forms of reuse.

\section{Architecture and acceptance}
Lean 4 combines an extensible elaborator with a kernel that checks proof terms
in dependent type theory~\cite{lean4}. Lean2LLM uses those interfaces to restrict
the role of the proposer. The proposer returns a tactic sequence. Lean executes
that sequence to construct a term for a proposition fixed by the caller.

\subsection{Fixed request and bounded proposal loop}
A request contains a name, a context, and a statement. The context and statement
remain fixed throughout a run. The native command \code{llm\_theorem} parses the
statement file as one complete Lean term and constructs the theorem syntax
structurally. This prevents a statement from closing one declaration and
introducing a different target through source concatenation. The native
\code{llm} tactic parses its input as a complete tactic sequence.

When no proposal file is supplied, the tactic emits a structured request with
Lean's actual goal. The external driver then invokes a proposer. Each retry
receives that goal, the fixed request, the previous candidate, Lean diagnostics,
an attempt number, and the remaining time. The driver bounds the attempt count,
individual Lean process time, source size, proof size, and memory, and checks an
elapsed search budget between calls. Toolchain discovery and module setup have
separate limits. A custom proposer is responsible for respecting its deadline;
the supplied CLI adapter enforces a subprocess deadline.

\begin{center}
\fbox{\begin{minipage}{0.92\linewidth}
\small
Fixed context and proposition $\longrightarrow$ Lean goal extraction
$\longrightarrow$ candidate proposal $\longrightarrow$ Lean elaboration and
kernel check $\longrightarrow$ axiom audit $\longrightarrow$ explicit-source
recheck $\longrightarrow$ stored proof.\par
Failed candidates return compiler diagnostics to the bounded proposal loop.
Retrieved candidates enter at the same checking stage.
\end{minipage}}
\end{center}

The result status \code{proved} means this run completed the configured
acceptance procedure. The status \code{unknown} covers exhausted attempts,
timeouts, missing toolchains, failed elaboration, rejected assumptions, and
proposer failures. It carries diagnostics and counters so that a user can
distinguish these cases. An \code{unknown} result establishes no negative
mathematical claim.

\subsection{Exact theorem and assumptions}
The command \code{audit\_llm} resolves the designated theorem in the
kernel-checked environment, verifies that the declaration is closed, and rejects
unresolved metavariables and proof placeholders. It traverses transitive axiom
dependencies. The inherited policy permits \code{propext},
\code{Classical.choice}, and \code{Quot.sound}. These assumptions are reported
as part of the checking profile.

The Python driver additionally requires a successful Lean process exit and the
expected audit marker. Acceptance also checks the generated explicit source
file, independently of the externally supplied tactic-file mechanism. This
second pass guards against discrepancies between the accepted candidate and
the exported theorem that users will receive. A lexical filter rejects a
conservative set of commands and metaprogramming constructs before execution.
The native axiom audit supplies a separate logical check: filtering text alone
cannot determine all assumptions in a proof's dependency closure.

\subsection{Execution boundary}
The inherited local runner uses macOS \code{sandbox-exec}. Checking processes
have no network permission or filesystem writes, and readable content is
restricted to the temporary inputs, pinned toolchain, trusted modules, and
selected runtime paths. This implementation fails closed when the required
isolation facility is unavailable. Its recorded tests include actual denied
canary reads and writes. Network restrictions are configured; the inherited
repository reports no independent penetration test of them.

\subsection{Browser workbench}
The web workbench uses a separate core Lean WebAssembly build in a fresh
isolated worker for each check. Its runtime wrapper is vendored from Browser Lean commit
\code{a8aed095}; the compiler identifies Lean commit \code{62b6a229}.
Runtime asset hashes and license notices accompany the source~\cite{browserlean}.
The initial compressed core payload is approximately 47.6 MiB. The deployed
profile excludes optional Mathlib and other extension layers. It bounds source
size and checking time; cancellation disposes of the worker. Browser HTTP
caching can reuse downloaded runtime assets across checks. Each submission
starts a new compiler worker, keeping the compiler's execution state isolated.
Local patches to the pinned runtime propagate compiler errors and aborts to
the application. A runtime fault permits one retry in another fresh worker
within the overall checking budget. Every completed or cancelled check
disposes of its worker.

The browser assembles a restricted theorem source and requires successful Lean
completion and a matching axiom report for the requested theorem. Its guards
reject placeholders, unsupported commands, and custom metaprogramming. The
browser profile has its own toolchain pin and checker fingerprint, separate
from the native bridge. A successful browser check is a local observation.
Records can be saved on the device, searched, and exchanged as JSON with the
Python library. Imported verification status is reset. A subsequent native
check resolves a native fingerprint and creates its corresponding content
identity. The site also supports an optional configured account-based storage
adapter. Remote server execution would require its own restricted worker and
resource policy; the current compiler executes in the browser.

\section{Two layers of proof reuse}
\subsection{Exact requests and a rechecked execution cache}
The original cache key hashes the normalized request together with a checker
fingerprint. That fingerprint includes the Lean version, executable bytes,
native bridge source and compiled module, and Python checking policy. Paths
also enter the original fingerprint, making it conservative and specific to the
local installation. A changed goal or checker identity selects a different
entry. This cache is appropriate for repeating the same experiment.

On a hit, the driver reads the stored request and proof, validates the key and
request, reruns Lean and the axiom audit, and rechecks the explicit export.
Damaged exports can be reconstructed from a freshly checked stored proof.
Consequently a cache hit saves proposal work; it still pays verification cost.
The cache is a performance aid under a fixed checking environment.

\subsection{Content-addressed records and indexed retrieval}
The proof library separates the identity of an artifact from its local
verification status. It stores records in SQLite and indexes title, statement, and tags
with full-text search. The record identifier is a SHA-256 digest of canonical
content. The shared JSON format requires valid Unicode and lossless integer
numbers within $\pm(2^{53}-1)$; fractional values and larger numbers travel as
strings. Python floats, including \code{1.0}, are rejected before hashing. The
schema version forms part of the exchange format. The content
contains the request, tactic body, exported source, backend name, toolchain
pin and fingerprint, tags, and provenance metadata.

The verification object is local evidence associated with that content. It
contains a status, scope, checking time, diagnostics, and checker fingerprint.
Import discards incoming verification claims. New content starts with an
unverified local status. When identical content already has a local result,
import preserves that existing evidence. Every subsequent reuse still requires
a fresh check under the receiving environment's selected profile.

\begin{table}[htbp]
\centering\small
\begin{tabular}{>{\raggedright\arraybackslash}p{0.21\linewidth}>{\raggedright\arraybackslash}p{0.70\linewidth}}
\toprule
Field & Purpose \\
\midrule
Schema and digest & Version the exchange format and identify canonical content.\\
Request & Preserve the exact context, proposition, and target name.\\
Proof and source & Retain both proposed tactics and an ordinary Lean export.\\
Toolchain & Record the pin and checker/dependency fingerprint.\\
Metadata & Carry title, tags, origin, licenses, and related artifact references.\\
Verification & Record who or what checked which scope, when, and with which profile.\\
\bottomrule
\end{tabular}
\caption{Logical content and verification evidence have separate roles in a
portable proof record. A content hash establishes byte-level identity under
the serialization rule; a Lean check establishes acceptance in an environment.}
\label{tab:record}
\end{table}

Search returns candidates. An exact replay query additionally matches the
statement, context, and toolchain identity. Broader queries may find a theorem
with useful quantified hypotheses or a lemma that can help prove another goal.
The retrieved theorem must be available in the target environment, and Lean
must check its actual application. Neither a shared title nor nearby embedding
coordinates establish applicability.

An artifact can be general in the ordinary mathematical sense: a theorem
quantified over a type, a function, and a natural number applies to every
instantiation satisfying its premises. Reusing that theorem requires checking
those premises. Automated discovery of useful substitutions, equivalences, or
isomorphisms can be layered above retrieval as a proposal mechanism. Its output
must include a checked application or transport proof. The storage index
provides no general complexity bound for that search.

\subsection{Portable bundles and identity}
A portable source bundle contains the proof record, its source file, an
integrity manifest, pinned bridge and Lake configuration, the checking driver,
and a verification entry point. After unpacking, a reader can run the entry
point with Python and the required local Lean toolchain without installing the
library package. Import validates digests, archive paths, sizes, and symbolic
link restrictions before storing the artifact. It does not execute bundled
code. The canonical digest covers the content being addressed; verification
metadata is excluded from the identity so that a later check can update local
status without pretending the mathematical content changed.

A complete dependency identity needs more than a display string such as
``Lean 4.'' The exact toolchain pin, the bridge and checker policy, and the
dependency lock manifest determine a verification profile. For richer projects,
the dependency closure should include package revisions, imported compiled
artifacts, relevant build options, and trusted plugins. A fingerprint is only as
complete as the inputs the implementation actually includes. The current
restricted Lean/Std bridge has a smaller dependency surface than a Mathlib
project; these profiles should be reported separately.

\subsection{Python interface}
The distribution is named \code{lean2llm}, with the import namespace
\code{lean\_llm}. Storage and search use Python's standard library, including
SQLite. Native checking additionally requires the pinned Lean environment and
the inherited macOS isolation facility. The following example stores supplied
tactics, checks them, searches the catalog, and exports a portable bundle:
\begin{lstlisting}
from lean_llm import ProofStore

with ProofStore("proofs.sqlite3") as store:
    record = store.add(
        {"name": "truth", "context": "", "statement": "True"},
        "exact True.intro", title="Truth example", tags=["example"])
    checked = store.check(record["id"])
    assert checked["check"]["accepted"]
    candidates = store.search("truth")
    store.export_bundle(checked["record"]["id"], "truth.zip")
\end{lstlisting}
Adding, importing, searching, and checking supplied proofs use no model calls.
The \code{solve} API takes an explicit proposer and an optional store. Its reuse
path searches an exact compatible request, rechecks the candidate, and returns
the resulting local record. A fresh proposal run uses the same acceptance path.

\section{Standalone certificates and statement alignment}
The public \code{navier-stokes-proof} repository provides a concrete portability
case study~\cite{sulea-certificates}. Its standalone package contains compressed
kernel declarations, a pinned challenge, a small checker based on Lean and Std,
licenses, provenance, hashes, and a verification entry point. The checker
compares designated theorem statements and their referenced definitions against
the supplied challenge, audits permitted assumptions, and replays the solution
declarations through Lean's kernel. Required mathematical dependencies travel
inside the certificate data, so replay avoids a separate Mathlib installation.

This profile differs from submitting a tactic sequence. The bundle pins Lean
4.34.0-rc2, while the initial bridge pins 4.33.1. The library therefore treats
certificate replay as an explicit backend with its own verifier identity. A
user supplies the trusted local verification program, its expected SHA-256
digest, and the bundle root; imported metadata cannot select an executable.
The adapter distinguishes \code{integrity-checked} from
\code{verifier-succeeded}; the latter records successful execution of the
explicit verifier, whose actual theorem-checking scope remains documented.
A source export and a kernel certificate
can share catalog and provenance conventions while retaining different check
procedures and resource budgets.

Formal acceptance also leaves a statement-alignment task. The certificate
repository includes a written comparison between the exact formal endpoints
and the corresponding informal Navier--Stokes alternatives. Its sources and
two AI review reports are preserved separately. The written translation is a
conventional argument, and its own status differs from a compiled theorem.
Lean2LLM adopts this separation in its export template: a package records the
formal statement, its assumptions, an informal description, the direction of
the claimed implication, unresolved obligations, and the scope of any review.
That structure helps a future reader distinguish a checked formal result from
the additional interpretation required to use it in another setting.

The reusable template follows five stages: preserve the source and licenses;
pin the formal target and dependencies; export the proof and recursive
dependencies; run the designated verifier and retain its output; document the
relationship between the formal target and the intended claim. A small source
bundle may require only the first, second, fourth, and fifth stages. The
certificate format serves cases where a standalone dependency closure is
available and worth its additional storage cost.

\section{Evidence, evaluation, and reproduction}
\subsection{Inherited public results}
Table~\ref{tab:evidence} distinguishes evidence published with the two
predecessor repositories. These are reported historical runs, inspected for
this manuscript. Their timings depend on the local machine, system load, and
toolchain. These records supply regression examples and provenance for the
implementation.

\begin{table}[htbp]
\centering\small
\begin{tabular}{>{\raggedright\arraybackslash}p{0.23\linewidth}>{\raggedright\arraybackslash}p{0.43\linewidth}>{\raggedright\arraybackslash}p{0.22\linewidth}}
\toprule
Artifact & Recorded result & Scope \\
\midrule
Initial bridge tests & 29 passing tests; 186.933 seconds in the committed log & Orchestration and real Lean checks\\
Induction demonstration & One live model proposal; three Lean checks; approximately 14 seconds & One quantified theorem\\
Standalone certificate & 94,719 solution declarations replayed; 580.518 seconds & Two pinned formal endpoints\\
\bottomrule
\end{tabular}
\caption{Inherited evidence, with distinct checking scopes. The bridge source
revision is \code{a409e42}; the inspected certificate repository revision is
\code{3b01518}. The certificate replay was recorded on 12 September 2026.}
\label{tab:evidence}
\end{table}

The bridge tests exercise accepted and rejected candidates, goal extraction,
compiler feedback before retry, time budgets, exact-cache reuse, corrupted
cache records, statement injection, direct transitive axiom auditing, and
sandbox denial of file operations. They use deterministic candidates with the
real Lean compiler. The live demonstration proves a recurrence property for an
arbitrary involution: if $f(f(x))=x$ for every $x$, then applying $f$ twice per
round returns $x$ for every natural-number round count. Induction proves that
statement. This result illustrates a reusable quantified theorem whose
hypothesis remains explicit.

The certificate run is evidence from a different repository and checker
profile. Its declaration count measures replayed proof data. It does not count
independently generated theorems, and it is unrelated to the bridge's proof
search success rate. The present project packages the storage and verification
workflow; the source repository attributes the underlying mathematical
formalization to its upstream authors.

\subsection{Evaluation of the library layer}
The installed \code{lean2llm==0.2.0} wheel passed 30 tests in a local
35.274-second run: 22 storage, certificate-adapter, and integrity tests, plus
eight native Lean integration tests. The run used a clean virtual environment,
Python's isolated execution mode, and a working directory outside the source
repository. It made zero model calls. The storage cases cover
deterministic identities, exact matching, full-text retrieval, bundle round
trips, corruption and unsafe archive rejection, and reset of imported
verification status. The native cases check fixed-statement acceptance,
rejection of unsuitable tactics and placeholders, transitive custom-axiom
rejection, and fresh native checking during reuse. They also check that a failed
recheck clears a prior local status. An end-to-end case imports a false proof
with a forged successful verification label, resets that label, and observes
rejection by the real Lean checker. A source bundle was unpacked and verified
using isolated Python execution outside the installed package. The additional
certificate case rejects a toolchain-file reference that escapes through a
symbolic link. These are new implementation results, separate from
Table~\ref{tab:evidence}.

The final audit also exercised the shared numeric and Unicode boundary and
rejected contradictory or wrongly targeted browser axiom reports. Browser
verification requires exactly one matching report for its private qualified
target. The frontend suite passed 19 tests; its scope includes source guards,
record import, checker-output interpretation, and mocked worker lifecycle,
runtime-fault recovery, and the one-retry bound.

A manual Chrome regression ran three checks in the same page without a reload:
the bundled paired-steps induction theorem was accepted, a tactic offering a
proof of \code{True} for the proposition \code{False} was rejected with a type
mismatch, and the restored induction theorem was accepted. The retained report
identifies these three inputs and their outcomes.

A separate interoperability run imported a browser-generated record, including
Unicode metadata, into Python. Its incoming successful status was reset to
\code{unverified}. Two real native Lean checks accepted the supplied proof and
produced a native verification record with a new toolchain fingerprint and
content identifier. A third report inspected the actual public standalone
Navier--Stokes template: nine manifest and referenced-file entries passed
integrity checking. That inspection did not execute the standalone verifier or
repeat the 94,719-declaration kernel replay.

The deterministic suite evaluates these mechanisms within its test cases.
Broad model accuracy requires a separately specified theorem set, budget,
baseline, and repeated-run protocol. Search ranking quality and large-catalog
latency remain subjects for further measurement. No comparative throughput or
proof-discovery success rate is inferred from the storage tests.

\subsection{Reproduction procedure}
A reproducible run begins from a pinned source revision and the version in
\code{lean-toolchain}. The inherited source-only regression commands are:
\begin{lstlisting}
lake build Llm
python3 -m unittest test_bridge -v
lake env lean examples/Verified.lean
\end{lstlisting}
The request appears in \path{examples/involution.json}; its accepted theorem is
in \path{examples/Verified.lean}. A reader can
check the supplied proof without invoking a model. A fresh model run requires
the configured proposer and its own credential mechanism; quotas and model
configuration belong in that run's provenance.

The installed-wheel regression log, run record, and companion interoperability
and integrity reports reside together:
\begin{lstlisting}
examples/public_platform/
  installed-wheel-tests.txt
  installed-wheel-result.json
  browser-interop-result.json
  navier-template-integrity-result.json
\end{lstlisting}
The test run used Python 3.13.3, the pinned Lean 4.33.1 distribution, and the
native-test opt-in flag. With the wheel installed in a clean environment and
the toolchain discoverable, the command has this form:
\begin{lstlisting}
LEAN_LLM_NATIVE_TESTS=1 \
  /path/to/venv/bin/python -I -m unittest discover \
  -s /path/to/repo/tests -v
\end{lstlisting}

Library reproduction additionally preserves the schema version, canonical
record, exported bundle, dependency profile, and local verification report.
The artifact identifier supports integrity checks and deduplication. The
verification report establishes what was run and the observed outcome. A
downloadable proof should provide both. The paper source and submission
metadata are included alongside the implementation so that changes to system
behavior can be reflected in the accompanying description.

\section{Related work and limitations}
LeanDojo exposes Lean proof environments and premise information; its ReProver
uses retrieval to select premises for generated proofs~\cite{leandojo}.
Lean Copilot integrates language-model inference into Lean to suggest proof
steps, complete goals, and retrieve premises~\cite{leancopilot}. These systems
establish that model proposals, Lean feedback, and retrieval are existing
approaches. Lean2LLM emphasizes a small externally orchestrated loop and the
subsequent artifact lifecycle: exact replay, indexed storage, portable exchange,
and explicit local rechecking. Its paper makes no priority claim for neural
theorem proving or retrieval-augmented proof search.

The initial proposer can fail on true theorems, repeatedly generate unsuitable
tactics, or spend its entire budget on an unproductive approach. Caching reduces
repeated proposal cost when an applicable proof already exists. A searchable
catalog enlarges the set of discoverable candidates. Neither mechanism yields
a completeness theorem or a polynomial bound for arbitrary proof search.

The Python driver, database, and web orchestration are covered by executable
tests and code review; they have no complete formal correctness proof. The
acceptance claim relies on Lean's checking implementation, the selected
toolchain and dependencies, the operating system, and correct association of a
checked theorem with the requested target. Hashes detect modification relative
to a trusted digest; an attacker who controls the reference digest can change
both content and manifest. Exporters, archive parsers, and certificate readers
remain software components requiring careful validation.

Axiom policies must be visible and versioned. A proof relying on a premise in
its statement establishes an implication under that premise. A user importing
a theorem also needs its definitions and scope. A successful formal check
does not supply missing empirical premises or an omitted correspondence to an
informal claim. These boundaries motivate preserving source, dependencies,
statement alignment, and verification evidence together.

\section{Conclusion}
Lean2LLM organizes a proof's lifecycle around a fixed statement and a repeatable
checking procedure. The original proposal loop produces Lean source under
explicit budgets and assumptions. Exact-cache replay rechecks that source.
The proof library makes artifacts addressable, searchable, exportable, and
eligible for a fresh local check. A certificate profile extends the same
cataloging principles to standalone kernel data. The resulting workflow gives
researchers a concrete way to keep, find, inspect, and reuse proofs after the
interaction that first produced them.

\paragraph{AI assistance and responsibility.}
OpenAI Codex assisted with software implementation, source inspection, test
construction, literature lookup, and drafting this manuscript. A language model
proposed the inherited induction proof; Lean checked the accepted proof term.
The certificate repository separately records its own AI review provenance.
Octavia Sulea is the author and retains responsibility for review, attribution,
and any decision to release or submit this manuscript. Automated checks and AI
reviews are described according to their actual scope.

\begin{thebibliography}{9}
\bibitem{lean4}
Leonardo de Moura and Sebastian Ullrich.
\newblock The Lean 4 Theorem Prover and Programming Language (System Description).
\newblock In \emph{Automated Deduction -- CADE 28}, 2021, pp. 625--635.
\newblock \url{https://doi.org/10.1007/978-3-030-79876-5_37}.

\bibitem{leandojo}
Kaiyu Yang, Aidan M. Swope, Alex Gu, Rahul Chalamala, Peiyang Song,
Shixing Yu, Saad Godil, Ryan Prenger, and Anima Anandkumar.
\newblock LeanDojo: Theorem Proving with Retrieval-Augmented Language Models.
\newblock \emph{Advances in Neural Information Processing Systems}, 36, 2023.
\newblock \url{https://arxiv.org/abs/2306.15626}.

\bibitem{leancopilot}
Peiyang Song, Kaiyu Yang, and Anima Anandkumar.
\newblock Lean Copilot: Large Language Models as Copilots for Theorem Proving in Lean.
\newblock arXiv:2404.12534, version 3, 2025.
\newblock \url{https://arxiv.org/abs/2404.12534v3}.

\bibitem{sulea-bridge}
Octavia Sulea.
\newblock Lean--LLM: local proposal, checking, and cache bridge.
\newblock Source, examples, and recorded test output, revision
\code{a409e42c806426f6f26ad76787fcebf902898d78}, 2026.
\newblock \url{https://github.com/mary-octavia/lean-llm/tree/a409e42c806426f6f26ad76787fcebf902898d78}.

\bibitem{sulea-certificates}
Octavia Sulea.
\newblock Navier--Stokes proof: standalone certificate checking and statement alignment.
\newblock Reproducibility artifacts, revision
\code{3b01518b6e33ab99b36a582219aa5c605be5268a}, 2026.
\newblock \url{https://github.com/mary-octavia/navier-stokes-proof/tree/3b01518b6e33ab99b36a582219aa5c605be5268a}.

\bibitem{browserlean}
Browser Lean contributors.
\newblock Browser Lean: a Lean WebAssembly runtime and browser integration.
\newblock Vendored source revision
\code{a8aed095534e67061e2626b94e8e69193099091d}; runtime provenance and asset
hashes preserved in the Lean2LLM source tree, 2026.
\newblock \url{https://github.com/live-codes/browser-lean/tree/a8aed095534e67061e2626b94e8e69193099091d}.
\end{thebibliography}
\end{document}
